\chapter*{A note to the reader}
\addcontentsline{toc}{chapter}{A note to the reader}
Ramanujan's tau function is easy to define and remarkably hard to keep away from zero. Its coefficients arise from one fixed infinite product,
\[
\Delta(q)=q\prod_{m\ge1}(1-q^m)^{24}=\sum_{n\ge1}\tau(n)q^n.
\]
The question in this volume is whether every one of those integer coefficients survives. Lehmer's conjecture says that it does. A great deal of structure is known, including strong upper bounds, exact recurrences, congruences and large verified ranges. None of these facts, by itself, provides a lower bound that excludes an exact zero at every prime.

This reader develops the Operation Komorebi packet program as a mathematical argument. We begin with the classical modular form, move through its convolution and five-core descriptions, construct permitted changes to its signed packet sums, and study what those changes preserve. The strongest conclusions concern the arithmetic of primitive packets and the exact limits of several proposed finishing mechanisms. They sharpen the problem rather than solve it.

\begin{claimboundary}
This edition does not prove Lehmer's conjecture and does not present an integer counterexample. It proves and explains auxiliary statements, records exact finite certificates and identifies the remaining obligation for the actual modular form $\Delta$. Written proofs, source-dependent theorems and bounded computations retain their distinct meanings throughout.
\end{claimboundary}

The primitive residual $R(n)$ plays a central role. It removes packets having a common coordinate factor. We will prove that $R(s^2)$ never vanishes and that $R(n)=0$ exactly when $\tau(n)=0$. These are useful structural theorems. They do not dispose of a prime zero, because $R(p)=\tau(p)$. At a prime the radial extraction leaves the original coefficient untouched.

The visual guide accompanies this book as a separate reading path. Its diagrams show cancellation, refinement, conserved content and the quantifier distinction behind the congruence argument. The source and evidence bundle contains the retained research packages, the edition's replay tools and the inputs required to inspect the finite claims. Its files are part of the mathematical record, but the reader need not follow an execution log to follow the proofs here.

This is Series 07, authored by Komorebi P. for Operation Komorebi. The reader edition is dated 9 October 2026. The underlying Lehmer research checkpoints were recorded on 26 September 2026 and have been preserved with their original scope. No claim of historical priority follows from their inclusion.

\chapter*{How to read the argument}
\addcontentsline{toc}{chapter}{How to read the argument}
The first five chapters establish the object and its exact residual descriptions. Chapters 6 and 7 turn the coefficient into a finite geometric sum and a genuine auxiliary polynomial. Chapters 8--10 explain why some polynomial corrections preserve every coefficient and show what a real sign certificate looks like. Chapters 11--14 distinguish geometric refinement from arithmetic conservation. Chapters 15 and 16 prove the primitive nonvanishing and zero-equivalence theorems. Chapters 17 and 18 extend the conservation law and separate interpolation from universal nullness. Chapters 19--21 examine actual congruence survivors, a compatible local trace-zero object and an arithmetic countermodel. The final chapter isolates the all-prime statement still needed.

Some prerequisites are used as contracts rather than reproved from the foundations. These include the modularity and Hecke properties of $\Delta$, harmonic theta modularity, the level-one valence formula, Deligne's representations and bound, Chebotarev's theorem, and the filtered-module classification. Their precise applications and normalizations are stated where used. The bibliography and source appendix identify the relevant primary or author-hosted passages. A reader may take those contracts as established and verify the project-specific algebra in the main text.

\begin{resultbox}[Three meanings of evidence]
A universal theorem has quantified hypotheses and a proof. A finite certificate proves a claim about its stated shell, matrix or polynomial by exact arithmetic. A bounded search describes its tested range and can locate examples or counterexamples within that range. A large count of passing tests does not change a bounded search into a universal theorem.
\end{resultbox}

The text uses $v_\ell(x)$ for the exponent of the prime $\ell$ in a nonzero integer $x$, with $v_\ell(0)=\infty$. The quadratic character $\chifive$ is zero on multiples of $5$, is $1$ on residues $1,4$, and is $-1$ on residues $2,3$. A shell is a set of integral vectors with a fixed norm and prescribed coordinate residues. A packet is one vector in that shell. Its weight can be positive or negative, but is never zero.

The recurring difficulty is cancellation. A representation as a sum of nonzero integers does not imply a nonzero sum. A method that changes individual contributions must state exactly why the total is unchanged and exactly why the changed contributions have a common nonzero sign. A method that supplies divisibility conditions must track whether its modulus is fixed or depends on the prime. These distinctions will do much of the conceptual work.
